\chapter{Implementation}
\label{ch:implementation}

% --------------------------------------------------------------------------
\section{Pure-Python BPF Encoder}
\label{sec:encoder}
% --------------------------------------------------------------------------

The language model outputs verifier-log assembly text; the kernel's
\texttt{BPF\_PROG\_LOAD} system call requires raw bytecode as a sequence of
8-byte little-endian instruction structs.
The pure-Python BPF encoder is the bridge between these two representations.
It is invoked by the reward function for every generated program before the
bytecode is passed to the \texttt{kcov\_validator} binary.

The BPF instruction struct layout is:
\begin{center}
  \texttt{opcode} (1 B) $\mid$
  \texttt{dst\_reg:4 | src\_reg:4} (1 B) $\mid$
  \texttt{off} (2 B, signed) $\mid$
  \texttt{imm} (4 B, signed)
\end{center}
totalling 8 bytes per instruction, packed little-endian.

The encoder processes each line of model output in three steps.
First, it normalises the line: the instruction index \texttt{N:}, the parenthesised
opcode byte \texttt{(XX)}, register-state annotations, and comment suffixes are
stripped, leaving the bare mnemonic text.
Second, it infers the instruction from that mnemonic. A priority-ordered set of
regular expressions, one per BPF instruction class (ALU, JMP, LD/ST, CALL, EXIT),
matches the mnemonic and looks the opcode up in a class-keyed table; the
destination register, source register (where present), 16-bit signed offset, and
32-bit signed immediate are read from the same match.
Third, it packs the opcode and the four operand fields into the 8-byte struct
using Python's \texttt{struct.pack} with the little-endian format specifier.
The opcode byte the verifier log prints is not consumed: it is stripped in the first
step and re-derived from the mnemonic in the second, so the encoder depends only on
the textual mnemonic, not on the log's opcode annotation.

Re-deriving the instruction from the mnemonic rather than invoking an assembler has
a mixed consequence for fuzzing. On one side the encoder applies no clang-style
validity check to operands: unconventional register assignments, out-of-range
immediates, and odd offsets within a recognised instruction form pass through
unchanged and reach the verifier as-is, where a clang-based assembler would reject or
normalise many of them. On the other side the encoder is itself a filter on the
opcode space --- it can only emit opcodes present in its per-class tables, so a
mnemonic outside those tables, or a line it cannot parse, is silently dropped rather
than passed through. The diversity the encoder preserves is therefore in the operands
of recognised instructions, not in arbitrary opcode bytes.

A line the encoder cannot parse is skipped, not fatal, so a single malformed line
does not discard an otherwise-valid program. An encode failure is declared only when
\emph{no} line parses: the encoder then returns no bytecode and the reward function
maps the empty result to a reward of~0.0, signalling the failure to GRPO without
crashing the training loop. When the parsed program does not already end in
\texttt{exit}, the encoder appends one so the bytecode is well-formed for loading.

The encoder is tested by 36 unit tests in
\texttt{tests/test\_reward\_encoder.py}.
The suite covers all six BPF instruction classes, boundary values for offset and
immediate fields (minimum and maximum signed 16- and 32-bit integers), programs
with zero instructions, programs with a single instruction of each class, and
malformed input strings of several types (missing opcode field, non-hex opcode,
truncated lines, empty string).

% --------------------------------------------------------------------------
\section{\texttt{kcov\_validator} Go Binary}
\label{sec:kcov-validator}
% --------------------------------------------------------------------------

The \texttt{kcov\_validator} binary (located in \texttt{tools/kcov\_validator/})
is the component that submits a BPF program to the kernel and returns both the
verifier verdict and the KCOV coverage trace for that program load attempt.
It runs inside the KCOV-instrumented virtual machine and is invoked remotely
over SSH by the reward function.

The binary's interface is minimal by design: it reads a hexadecimal-encoded
bytecode string from standard input and writes a single JSON object to standard
output:

\begin{lstlisting}[caption={\texttt{kcov\_validator} output format}]
{"verdict": "ACCEPTED", "pcs": [18446744071584860160, ...]}
{"verdict": "REJECTED", "pcs": [18446744071584860160, ...]}
\end{lstlisting}

Internally, the binary executes six steps in sequence.
\begin{enumerate}
  \item Open the KCOV device file \texttt{/sys/kernel/debug/kcov} and
        \texttt{mmap} the ring buffer into the process address space.
  \item Call \texttt{ioctl(KCOV\_ENABLE)} with trace mode
        \texttt{KCOV\_TRACE\_PC} to begin recording PCs for the current task.
  \item Decode the hex input to raw bytes; invoke the \texttt{BPF\_PROG\_LOAD}
        system call with a minimal program attributes struct.
        The verifier verdict (accepted or rejected) and the error log are
        captured from the system call return value.
  \item Call \texttt{ioctl(KCOV\_DISABLE)} to stop recording.
  \item Call \texttt{readPCs()}: read the count $n$ from element~0 of the
        ring buffer, then collect elements~$1, \ldots, n$ into a slice of
        \texttt{uint64} values.
  \item Serialise \texttt{\{verdict, pcs\}} to JSON and write to stdout.
\end{enumerate}

\textbf{Unconditional trace collection.}
KCOV trace collection is unconditional: both ACCEPTED and REJECTED loads return
the full PC array covering the kernel code paths executed during the load
attempt.
This is a requirement of the reward signal, not an incidental choice.
A verdict-conditional variant --- one that returned an empty trace for rejected
programs --- would make every rejected program indistinguishable to the reward
function, since they would all carry the same empty coverage.
That homogeneity collapses the within-group reward variance GRPO depends on; it
is the failure mode analysed in Section~\ref{sec:rl-plateau}.

\textbf{Coverage scope.}
The trace records every kernel program counter executed between
\texttt{KCOV\_ENABLE} and \texttt{KCOV\_DISABLE}, which bracket the whole
\texttt{BPF\_PROG\_LOAD} system call rather than the verifier alone. It therefore
includes the surrounding load-path machinery --- argument copy-in, program and map
allocation, and, for accepted programs, JIT compilation and teardown --- and no
address-range filter restricts it to the verifier translation unit, in contrast to
the per-region KCOV filtering syzkaller applies. The PC counts reported in this
thesis are accordingly \texttt{BPF\_PROG\_LOAD}-path coverage, not verifier-only
coverage, with two consequences. An accepted program traverses post-verification
paths that a rejected one never reaches, so the two verdict classes are not measured
over the same code region. And because part of the load path is exercised by every
program, the absolute counts carry a roughly constant load-path component on top of
the verifier paths that actually vary between programs.
Section~\ref{sec:diversity-wall} states how this bounds the central measurement.

The same scope reaches the RL reward, not only the offline measurement. The novelty
terms of Section~\ref{sec:reward-redesign} credit any rarely-visited PC, so a program
can earn novelty for background-kernel addresses --- allocation, soft-irq, and
scheduler paths --- that the noise characterization of
Section~\ref{sec:diversity-wall} shows vary independently of the program. The
exploration signal is contaminated to that extent. Restricting the trace to the
verifier translation unit with a per-region address filter would remove the
contamination and is left to future work (Section~\ref{sec:future-work}).

\textbf{VM configuration.}
The virtual machine runs Debian trixie with a Linux 6.8.0 kernel compiled with
\texttt{CONFIG\_KCOV}, \texttt{CONFIG\_KASAN}, and \texttt{CONFIG\_UBSAN} enabled.
It is launched by QEMU with a virtio-9p filesystem share that exposes the corpus
directory to the host.
The \texttt{CONFIG\_KASAN} and \texttt{CONFIG\_UBSAN} instrumentation layers are
retained from the data-collection phase; they have no effect on the KCOV trace
but provide crash-detection coverage during RL training.

% --------------------------------------------------------------------------
\section{GRPO Training Loop}
\label{sec:grpo-loop}
% --------------------------------------------------------------------------

The GRPO training loop is implemented in \texttt{ml/rl\_grpo.py} using the
\texttt{GRPOTrainer} class from TRL~0.14.0, pinned below version~0.15 for API
stability.
At each step the trainer samples $G$ completions per prompt from the current
policy and scores each one with the KCOV reward function of
Section~\ref{sec:reward-redesign}: the completion is encoded to bytecode,
submitted to \texttt{kcov\_validator} over SSH, and converted to a scalar reward
from the returned verdict and PC trace.
GRPO then forms group-normalised advantages from the $G$ rewards and applies the
policy-gradient update.
The reward function is the loop's only coupling to the kernel; every other part
of the step is standard policy-gradient training.

Making this loop run reliably required two library-level workarounds --- a
TRL/vLLM availability check and a Transformers KV-cache regression --- and long
unattended runs are guarded by run-monitoring and checkpoint-resume tooling.
These are reproduction details rather than part of the method; they are recorded
in Appendix~\ref{app:eng-notes}.

% --------------------------------------------------------------------------
\section{Remote Reward Pipeline}
\label{sec:remote-reward}
% --------------------------------------------------------------------------

The KCOV reward oracle cannot run in the cloud: \texttt{kcov\_validator} needs a
custom kernel built with \texttt{CONFIG\_KCOV} and raw
\texttt{/sys/kernel/debug/kcov} access, which cloud hypervisors do not expose.
Training, on the other hand, benefits from a cloud GPU with more VRAM than the
local 8.59~GB card. The pipeline therefore decouples the two: generation runs on a
cloud GPU while reward evaluation stays on the local KCOV VM, connected by an HTTP
bridge.

\texttt{tools/reward\_server.py} is a FastAPI server wrapping \texttt{ml/reward.py}.
It accepts a batch of completions, invokes the reward function --- which calls
\texttt{kcov\_validator} over SSH --- for each, and returns the scalar rewards. A
pre-shared API key guards the endpoint, and the trainer reaches the server over an
ngrok tunnel passed via the \texttt{--remote-reward-url} flag. If the server stays
unreachable the trainer aborts rather than training on missing rewards, so a
network failure cannot silently corrupt the policy.

This pipeline carried RL run~2 (Chapter~\ref{ch:results}).
